\documentclass[12pt]{article}
\usepackage[utf8]{inputenc}
\usepackage[T1]{fontenc}
\usepackage[margin=1in]{geometry}
\usepackage{setspace}
\usepackage{xcolor}
\usepackage[colorlinks=true, urlcolor=blue!60!black, linkcolor=blue!60!black]{hyperref}
\setlength{\parindent}{2em}
\setlength{\parskip}{1.25em}
\renewcommand{\baselinestretch}{1.0}
\begin{document}
\begin{center}
{\LARGE \textit{Plyler v. Doe}, 457 U.S. 202 (1982)\par}
\vspace{0.8em}
{\large Case Brief --- Four-Page Edition\par}
\end{center}
\vspace{1em}

\section{Memory Jogger}

By a 5--4 vote, the Court held that Texas could not deny a free public education to undocumented children: they are ``persons within its jurisdiction'' under the Equal Protection Clause, and a law imposing a lifetime hardship on innocent children must further some substantial goal of the State (\href{https://supreme.justia.com/cases/federal/us/457/202/}{457 U.S. 202}).

\section{Detailed Case Facts}

In May 1975, Texas amended Education Code \S~21.031 to withhold state funds from local districts for educating children not ``legally admitted'' into the United States and to let districts deny them free enrollment; only ``citizens of the United States or legally admitted aliens'' were entitled to tuition-free schooling. In July 1977, the Tyler Independent School District began requiring undocumented children to pay a ``full tuition fee'' of \$1,000 a year.

\section{Procedural History}

Undocumented children of Mexican origin in Smith County brought a class action against Superintendent James Plyler; Texas intervened. After an evidentiary hearing, the district court found that the law neither deterred illegal immigration nor improved education and enjoined it under the Equal Protection and Supremacy Clauses; the Fifth Circuit affirmed. In consolidated statewide litigation, \textit{In re Alien Children Education Litigation}, the Southern District of Texas applied strict scrutiny and invalidated the law, and the Fifth Circuit summarily affirmed. The Supreme Court heard both cases together on December 1, 1981, and decided them on June 15, 1982, without reaching preemption.

\section{Judicial Votes}

Affirmed, \textbf{5--4}.
\begin{itemize}
\item \textbf{Majority:} Justice Brennan, joined by Justices Marshall, Blackmun, Powell, and Stevens.
\item \textbf{Concurrences:} Justices Marshall, Blackmun, and Powell, each separately.
\item \textbf{Dissent:} Chief Justice Burger, joined by Justices White, Rehnquist, and O'Connor.
\end{itemize}

\section{Holding}

(1) Unanimously, undocumented immigrants physically present in a State are ``persons within its jurisdiction'' protected by the Equal Protection Clause. (2) By 5--4, although undocumented aliens are not a suspect class and education is not a fundamental right, \S~21.031 imposes a lifetime hardship on innocent children and cannot stand unless it furthers a substantial goal of the State, which Texas failed to show.

\section{Analysis of Opinions}

\textbf{Majority (Brennan).} An alien is a ``person'' under \textit{Yick Wo} and \textit{Wong Wing}, and the Reconstruction debates show that ``within its jurisdiction'' is territorial. Strict scrutiny does not apply, but ordinary rational basis review does not fit either: children ``can affect neither their parents' conduct nor their own status,'' and punishing them risks a permanent ``shadow population.'' Education, under \textit{Brown}, \textit{Yoder}, and \textit{Meyer}, sustains democratic institutions, and its denial inflicts the stigma of illiteracy. No federal policy supported exclusion, unlike the labor rule in \textit{DeCanas v. Bica}. Texas's interests failed: tuition for children was a ``ludicrously ineffectual attempt to stem the tide of illegal immigration,'' exclusion would not improve quality, and many children would remain in Texas.

\textbf{Concurrences.} Marshall reiterated his \textit{Rodriguez} sliding-scale view. Blackmun likened the total denial of education to denial of the vote. Powell, author of \textit{Rodriguez}, invoked the illegitimacy cases and a ``fair and substantial relation'' test, while faulting Congress's lack of ``effective leadership'' on immigration.

\textbf{Burger (dissenting).} Denying children schooling is ``senseless'' as policy, but the Court acted as ``Platonic Guardians.'' Once it conceded no suspect class or fundamental right, rational basis applied and Texas's fiscal interest sufficed; the ruling would ``likely stand for little beyond the results in these particular cases.''

\section{Examples: Future Applications}

\begin{enumerate}
\item \textbf{Same-Side (Unconstitutional): Proof of Status to Enroll.} Requiring documents showing lawful presence before enrollment produces the absolute, status-based deprivation \textit{Plyler} condemned.
\item \textbf{Same-Side (Unconstitutional): Tuition for Citizen Children.} Charging tuition to U.S. citizen children of undocumented parents punishes children for their parents' conduct.
\item \textbf{Opposite-Side (Likely Upheld): In-State University Tuition.} Denying in-state college tuition to undocumented graduates likely survives; K--12 tuition for unaccompanied minors raises the question in Powell's footnote.
\item \textbf{Opposite-Side (Upheld): Residency Rules.} A neutral bona fide residency requirement is valid, as \textit{Martinez v. Bynum} (1983) held, though a State may not relabel undocumented children as nonresidents.
\item \textbf{Fence-Sitter: Reporting Status.} A rule requiring schools to verify and report students' status, while keeping enrollment free, turns on proof of a chilling effect.
\end{enumerate}

\section{Critique}

Critics fault \textit{Plyler} for combining two non-triggers, a class that is not suspect and a right that is not fundamental, into ad hoc intermediate scrutiny, and argue that its logic fits Supremacy Clause preemption better than equal protection. Its limits show in contrast: \textit{San Antonio ISD v. Rodriguez} upheld funding disparities because every child still received a basic education, and \textit{Kadrmas v. Dickinson Public Schools} (1988) upheld school-bus fees, stressing \textit{Plyler}'s ``unique circumstances.''

\section{Key Quotations}

\begin{quote}
``Section 21.031 imposes a lifetime hardship on a discrete class of children not accountable for their disabling status. The stigma of illiteracy will mark them for the rest of their lives. By denying these children a basic education, we deny them the ability to live within the structure of our civic institutions, and foreclose any realistic possibility that they will contribute in even the smallest way to the progress of our Nation. In determining the rationality of \S~21.031, we may appropriately take into account its costs to the Nation and to the innocent children who are its victims. In light of these countervailing costs, the discrimination contained in \S~21.031 can hardly be considered rational unless it furthers some substantial goal of the State.''
\end{quote}
\noindent --- Justice Brennan, Opinion of the Court, 457 U.S. at 223--24.

\begin{quote}
``Public education is not a `right' granted to individuals by the Constitution. \textit{San Antonio Independent School Dist.\ v.\ Rodriguez}, 411 U.S. 1, 35 (1973). But neither is it merely some governmental `benefit' indistinguishable from other forms of social welfare legislation. Both the importance of education in maintaining our basic institutions, and the lasting impact of its deprivation on the life of the child, mark the distinction.''
\end{quote}
\noindent --- Justice Brennan, Opinion of the Court, 457 U.S. at 221.

\begin{quote}
``It is difficult to understand precisely what the State hopes to achieve by promoting the creation and perpetuation of a subclass of illiterates within our boundaries, surely adding to the problems and costs of unemployment, welfare, and crime. It is thus clear that whatever savings might be achieved by denying these children an education, they are wholly insubstantial in light of the costs involved to these children, the State, and the Nation.''
\end{quote}
\noindent --- Justice Brennan, Opinion of the Court, 457 U.S. at 230.

\end{document}
